% ---------------------------------------------------------------------------
% content/hebrewpages.tex -- the commands available in hebrew_pages.tex.
% ---------------------------------------------------------------------------

\linkedsection{rtl}{Right-to-Left Text and Latin Fragments}

\texttt{hebrew\_pages.tex} sets Hebrew as its default language with the \texttt{polyglossia} package, so ordinary paragraphs are typeset right to left and headings are aligned to the right.
Inline mathematics such as \verb|$f(k)\cdot n^{O(1)}$| works unchanged.
A Latin word or phrase inside a Hebrew sentence, however, must be marked, or the bidirectional algorithm will reorder it. The template defines

\begin{lstlisting}
\newcommand{\en}[1]{\beginL{\fontspec{Latin Modern Roman}#1}\endL}
\end{lstlisting}

so that \verb|\en{NP}| or \verb|\en{(fixed-parameter tractable)}| is typeset left to right, in a Latin font.
The year on the Hebrew cover page is set the same way, with \verb|\en{\thesisyear}|.

\linkedsection{hcite}{Citations: the \texttt{\textbackslash hcite} Command}
\label{sec:hcite}

A citation inside Hebrew text is a Latin fragment too, with one more requirement: it must be a link into the English bibliography.
The template's command is

\begin{lstlisting}
\newcommand{\hcite}[1]{\mbox{\LR{{\fontspec{Latin Modern Roman}\cite{#1}}}}}
\end{lstlisting}

and it is used exactly like \verb|\cite|: \verb|\hcite{Karp72,GJ79}|.
Three things make it work.

\begin{itemize}
    \item \verb|\LR| (from the \texttt{bidi} package) and the Latin font keep the citation label readable and left to right; \verb|\mbox| keeps it on one line.
    \item \texttt{hebrew\_pages.tex} loads the same \texttt{.bib} file with the same \texttt{biblatex} options as \texttt{main.tex}. The \texttt{alphabetic} style derives its labels from author names and years, so the Hebrew document prints \texttt{[Kar72]} for the same entry the English bibliography lists as \texttt{[Kar72]}, even though the Hebrew document prints no bibliography of its own.
    \item \texttt{biblatex} makes each citation a link to the anchor \texttt{cite.0@<key>}. In the Hebrew PDF alone that anchor is missing; after the merge it is found in the English bibliography and the link is rewired to it (\Cref{sec:merge}).
\end{itemize}

One rule follows from this: an entry can only be linked to if it appears in the English bibliography, that is, if it is cited (or listed with \verb|\nocite|) somewhere in \texttt{main.tex}. An entry cited only from Hebrew still prints its label, but the merge step reports its link as unresolved.

If you switch the English bibliography to a numeric style, switch the Hebrew document too, and keep in mind that numeric labels then depend on citation order and will differ between the two documents; the links still work, the printed numbers will not match. The alphabetic style avoids the problem.

\linkedsection{enref}{References to English Results: the \texttt{\textbackslash enref} Command}

To refer to an English theorem, figure, definition or table from Hebrew text, use

\begin{lstlisting}
\newcommand{\enref}[2]{\hyperlink{target:#1}{\englishTOCnumber{#2}}}
% usage, inside Hebrew text:  (<Hebrew word for Theorem> \enref{thm:main}{thm:main})
\end{lstlisting}

The first argument is the anchor name without its \texttt{target:} prefix, the second is the label.
The command prints the object's number, read from the English build, and makes it a link to the anchor, so the Hebrew text needs no hard-coded numbers and does not break when the English part is renumbered.
The Hebrew summary at the end of this document refers to \Cref{thm:main}, \Cref{def:anchor}, \Cref{fig:pipeline} and \Cref{tab:naming} this way.

Behind \verb|\enref| are the two macros that also drive the Hebrew table of contents:

\begin{lstlisting}
\newcommand{\englishTOCpage}[1]{\LR{\pageref*{en-#1}}}
\newcommand{\englishTOCnumber}[1]{\LR{\ref*{en-#1}}}
\end{lstlisting}

They print a page or a number from an \texttt{en-} label, left to right, and deliberately without a link of their own (the starred forms): the enclosing \verb|\hyperlink| is the link, and the merge script expects exactly one destination per link.

\linkedsection{hebrewtoc}{The Hebrew Table of Contents}
\label{sec:hebrewtoc}

The Hebrew TOC is written by hand, one line per English chapter or section, in \texttt{hebrew\_pages.tex}.
Each line combines the three ingredients above: a number, a linked Hebrew title and a page.

\begin{lstlisting}
\englishTOCnumber{chapter:build}.\ \hyperlink{chapter:build}{<Hebrew chapter title>}
\dotfill \englishTOCpage{chapter:build}\\[0.2em]
\hspace{2em}\englishTOCnumber{toc:sec:script}\ \hyperlink{target:sec:script}{<Hebrew section title>}
\dotfill \englishTOCpage{toc:sec:script}\\[0.1em]
\end{lstlisting}

Unnumbered chapters (the conclusion, the bibliography) omit the number.
The last two lines, for the Hebrew summary and abstract, use local labels, since those pages are in the Hebrew document itself; the English TOC reads the same labels through the \texttt{he-} prefix to list them on its side.
The order of the lines, and the Hebrew titles, are yours to maintain; the numbers and pages are not.
